\MathBookSourcePage{361}
\subsection{习题}
\label{sec:ch12-exercises}
\setcounter{exerciseitem}{0}

\begin{MBExercise}
\label{exr:ch12-coercivity-stokes-a}
证明~\eqref{eq:ch12-stokes-a-form} 中的形式 $a(\cdot,\cdot)$ 在 $\mathring H^1(\Omega)^n$ 上强制. 提示: 使用命题~\ref{prop:ch05-mixed-boundary-coercivity}.
\end{MBExercise}

\begin{MBExercise}
\label{exr:ch12-hdiv-hilbert}
证明 $H(\operatorname{div};\Omega)$ 关于内积
\[
\MathBookFormulaID{formula-ch12-ex02-hdiv-inner-product}
(u,v)_{H(\operatorname{div};\Omega)}
:=(u,v)_{L^2(\Omega)^n}
+(\operatorname{div}u,\operatorname{div}v)_{L^2(\Omega)}
\]
是 Hilbert 空间. 提示: 参照定理~\ref{thm:ch01-sobolev-banach} 的证明.
\end{MBExercise}

\begin{MBExercise}
\label{exr:ch12-symmetric-gradient-stokes-form}
对 Stokes 问题定义
\[
\MathBookFormulaID{formula-ch12-ex03-symmetric-gradient-form}
a(u,v):=2\int_\Omega\sum_{i,j=1}^n
\varepsilon_{ij}(u)\varepsilon_{ij}(v)\,dx,
\qquad
\varepsilon_{ij}(u):=\frac12(u_{i,j}+u_{j,i}).
\]
证明当 $u,v\in Z$ 时, 该形式等于~\eqref{eq:ch12-stokes-a-form}. 提示: 使用 $\sum_i u_{i,i}=0$.
\end{MBExercise}

\begin{MBExercise}
\label{exr:ch12-korn-variant-coercivity}
证明习题~\ref{exr:ch12-symmetric-gradient-stokes-form} 中的形式在 $Z$ 上强制. 提示: 使用第~\ref{chap:ch11-planar-elasticity} 章的 Korn 不等式.
\end{MBExercise}

\begin{MBExercise}
\label{exr:ch12-nonconforming-symmetric-gradient-not-coercive}
证明习题~\ref{exr:ch12-symmetric-gradient-stokes-form} 中形式的非协调版本在~\eqref{eq:ch11-nonconforming-vector-space} 定义的 $V_h^*$ 上不以与 $h$ 无关的常数强制(参见定理~\ref{thm:ch12-nonconforming-stokes-velocity}). 提示: 考虑 $\mathbb R^2$ 上的正则网格, 寻找具有重复模式的 $u$, 使 $a(u,v)=0$ 对所有具有紧支集的 $v\in V_h^*$ 成立.
\end{MBExercise}

\begin{MBExercise}
\label{exr:ch12-symmetric-gradient-mixed-formulation}
证明第~\ref{sec:ch12-abstract-mixed-formulation} 节的结果适用于习题~\ref{exr:ch12-symmetric-gradient-stokes-form} 中的形式. 提示: 使用习题~\ref{exr:ch12-korn-variant-coercivity}.
\end{MBExercise}

\begin{MBExercise}
\label{exr:ch12-finite-dimensional-inf-sup-nullspace}
设 $\Pi_h$ 是有限维空间. 证明若~\eqref{eq:ch12-discrete-inf-sup} 中 $\beta=0$, 则存在非零 $q\in\Pi_h$, 使得
$b(v,q)=0$ 对所有 $v\in V_h$ 成立. 提示: 在 $\Pi_h$ 的单位球面上使用紧性.
\end{MBExercise}

\begin{MBExercise}
\label{exr:ch12-operator-M-bound}
证明~\eqref{eq:ch12-operator-M-definition} 定义的算子 $M$ 有界, 且
\[
\MathBookFormulaID{formula-ch12-ex08-operator-M-bound}
\|M\|\leq\frac1\beta.
\]
提示: 在~\eqref{eq:ch12-discrete-inf-sup-pointwise} 中取 $q=Mu$, 并使用~\eqref{eq:ch12-operator-M-definition}.
\end{MBExercise}

\begin{MBExercise}
\label{exr:ch12-minimum-norm-Lp}
证明~\eqref{eq:ch12-minimum-norm-Lp}. 提示: 若 $v\in Z_h^p$, 则 $Lp-v\in Z_h$, 并且 $(Lp,Lp-v)_V=0$.
\end{MBExercise}

\begin{MBExercise}
\label{exr:ch12-operator-equivalence-converse}
证明若~\eqref{eq:ch12-operator-L-definition} 中的算子 $L$ 定义良好, 则~\eqref{eq:ch12-discrete-inf-sup} 以 $\beta=1/\|L\|_{\Pi\to V}$ 成立. 提示: 在~\eqref{eq:ch12-discrete-inf-sup} 中取 $v=Lq$.
\end{MBExercise}

\MathBookSourcePage{362}
\begin{MBExercise}
\label{exr:ch12-curl-Morley-inclusion}
证明图~\MBUnresolvedReference{fig:ch10-source-10-5}{10.5} 所示 Morley 空间的旋度属于定理~\ref{thm:ch12-nonconforming-stokes-velocity} 中使用的分片线性非协调空间. 提示: 跨一条公共边的切向导数跳跃是二次多项式, 在边的两个端点为零; 因而它在边中点也为零.
\end{MBExercise}

\begin{MBExercise}
\label{exr:ch12-full-mixed-constant-bound}
用 $C$, $\alpha$ 和 $\beta$ 给出推论~\ref{cor:ch12-full-mixed-quasioptimality} 中常数 $c$ 的一个界.
\end{MBExercise}

\begin{MBExercise}
\label{exr:ch12-full-mixed-strang}
证明~\eqref{eq:ch12-full-mixed-Strang}.
\end{MBExercise}

\begin{MBExercise}
\label{exr:ch12-nonconforming-pressure-error}
证明定理~\ref{thm:ch12-nonconforming-pressure-error}.
\end{MBExercise}

\begin{MBExercise}
\label{exr:ch12-local-test-function-norm}
证明~\eqref{eq:ch12-local-test-function-norm}.
\end{MBExercise}
